








Porosity ($p$) is the fraction of a dust grain's volume that is empty space. A low porosity corresponds to a more compact grain (like a rock), while a high porosity corresponds to a grain with more empty space (like a sponge).

Instead of specifying the porosity directly, the filling factor ($f$) is used, defined as:
\begin{equation}
 f = 1 - p,
    \label{eq: filling factor}
\end{equation}
\noindent where $f=1$ corresponds to a compact grain ($p$ = 0), and an $f$ value less than 1 corresponds to a more porous grain \citep{Zhang2023}. A range of filling factors was tested ($f = 1, 0.75, 0.5, 0.3, 0.25, 0.1, 0.01$) to sample different dust grain porosities. The results for $f = 1, 0.5, 0.3, 0.1$ are discussed in more detail, and some results for $f = 0.75, 0.25, 0.01$ are included.

% and the main four values are presented in Chapter \ref{sec: chapter5_spherical_dust_model_results}: $f = 1, 0.5, 0.3, 0.1$. 
%Four different filling factors were considered ($f = 1, 0.75, 0.5, 0.25$) to sample a wide range of dust grain porosities. 

For each filling factor, the \texttt{DSHARP} code was used to calculate the corresponding density and optical constants of the dust mixture. The resulting density values are shown in Table \ref{tab:dsharp_density}. As expected, the density decreases as the filling factor decreases and the porosity increases.  

\question{Add more detail on how the \texttt{DSHARP} code incorperates different filling factors?} \SC{A: You haven't defined opacity yet.  You may want to wait until you get to equations that use the porosity.}


\note{specify that irregular is compact}
%\add{Question from disks on the exe: different function? How does the code integrate the different porosities}

\begin{table}[h!]
\centering
\caption{Summary densities calculated using \texttt{DSHARP} for each filling factor.}
\begin{tabular}{ll}
\toprule
\textbf{Filling Factor (\boldmath$f$)} &
\textbf{Density (g/cm\boldmath$^3$)} \\
\midrule
1.0    & 1.675 \\
0.5 & 0.838 \\
0.3  & 0.503 \\
0.1 & 0.168 \\
\bottomrule
\end{tabular}
\label{tab:dsharp_density}
\end{table}

% \draft{With the density for each filling factor, and the radius of the dust grain, we can calculate the mass of the dust grain for a spherical grain with:
% \begin{equation}
%     \text{mass} = \rho \cdot \frac{4\pi}{3} a^3.
% \end{equation}
% }

% ----------------------------------------------
% ----------------------------------------------
















%
%
%
%
%

